% !TEX root = ../main.tex
\begin{abstract}
% CS197 move 1: problem motivation (failure mode) -- non-response has several explanations.
Non-response to cognitive behavioral therapy (CBT) does not identify a single failure of treatment, theory or delivery.
% CS197 move 2: assumption (setup bit) -- package efficacy does not settle component causality.
Evidence that a treatment package works does not establish which components cause improvement or why some patients do not benefit.
% CS197 move 3: contribution (bit flip) -- a critical synthesis, not a new systematic review.
We critically evaluate ten proposed reasons CBT may fail, using selected meta-analyses, component trials and mechanistic studies rather than claiming systematic coverage.
% CS197 move 4a: instantiation (bit flip instantiation) -- retain all ten reasons.
The review covers efficacy, mechanisms, patient and context factors, and delivery.
% CS197 move 4b: instantiation (bit flip instantiation) -- keep the original case study within the broader scope.
Executive function and adult attention-deficit/hyperactivity disorder provide the deepest case study.
% CS197 move 5a: evaluation (diagnostic result) -- evidence map; no newly pooled effect estimate.
CBT remains effective on average, but benefits depend on the comparator, and many limitations are shared with other treatments.
% CS197 move 5b: evaluation (diagnostic result) -- component results vary and nulls are not equivalence.
Component findings are mixed and often underpowered; they establish neither the universal necessity nor the irrelevance of cognitive restructuring.
For example, a depression review (16 studies, 22 comparisons) favoured full treatments over versions with one component removed ($g=0.21$, 95\% confidence interval $[0.03,0.39]$), despite substantial methodological limitations.
% CS197 move 5c: evaluation (diagnostic result) -- evolutionary inference and selection boundaries.
Evolutionary conservation motivates learning-based hypotheses but cannot rank therapies or establish that experience must precede cognition.
Researcher allegiance and publication selection warrant scrutiny without establishing CBT-specific reviewer suppression.
Acceptance, alliance, reactance and feedback research further qualify accounts that isolate technique from goals and delivery.
% CS197 move 6: implication (north star) -- proposed tests, not validated treatment-selection rules.
We propose falsifiable component and moderator tests of which interventions produce durable functional change, for whom and under what conditions.
\end{abstract}
